\section{The limiting Jacobi roots}
For each fixed $d$, the monic polynomials
\[
 L_j(x)=i^{-j}G_j(ix,R)
\]
obey
\[
 L_{j+1}(x)=xL_j(x)-j(j+R)L_{j-1}(x).
\]
Thus $L_k$ is the characteristic polynomial of the $k\times k$ real symmetric Jacobi matrix $J_{d,k}$ with zero diagonal and off-diagonal entries $\sqrt{j(j+R)}$, $1\le j<k$. Its roots $\lambda$ are real, simple, and symmetric. For every positive root put $\alpha=\lambda^2/4$. Then
\begin{equation}\label{aud:fr:rootproduct}
 G_k(2u,R)=(2u)^k\prod_{\lambda>0}(1+\alpha/u^2).
\end{equation}
Empty products are one, and empty root maxima are zero. The elementary matrix norm bound gives, for $k\ge2$,
\begin{equation}\label{aud:fr:rootbound}
 \max\alpha\le(k-1)(k-1+R).
\end{equation}

\begin{lemma}[A finite-mass arcsine mixture]\label{aud:fr:measurelemma}
If $k/d\to\kappa\ge0$, then the finite measures
\[
 \nu_{d,k}=\frac1d\sum_{L_k(\lambda)=0}\delta_{\lambda/d}
\]
converge weakly to the measure characterized by
\begin{equation}\label{aud:fr:measure}
 \int f\,d\nu_\kappa
 =\frac1\pi\int_0^\kappa\int_0^\pi
 f\bigl(2\sqrt{x(x+4)}\cos\theta\bigr)\,d\theta\,dx.
\end{equation}
The measure has total mass $\kappa$. For each continuous $f$ on a common compact spectral interval, convergence to the right-hand side with $\kappa=k/d$ is uniform when $k/d$ ranges over a compact interval.
\end{lemma}
\begin{proof}
The norm bound gives a common compact spectral interval whenever $k/d$ is bounded. Odd trace moments vanish. For a fixed even power $2p$, expand the diagonal entries of $(J_{d,k}/d)^{2p}$ as sums over nearest-neighbor returning paths. Away from the $O(p)$ boundary rows, there are $\binom{2p}{p}$ such paths. At a row with $j/d\to x$, every path product tends to $[x(x+4)]^p$. The convergence is uniform on compact $x$ intervals, because $\sqrt{x(x+4)}$ is continuous there. The boundary rows contribute $o(1)$ after dividing the trace by $d$. Riemann sums therefore give
\[
 \frac1d\operatorname{tr}(J_{d,k}/d)^{2p}
 \longrightarrow\binom{2p}{p}\int_0^\kappa[x(x+4)]^p\,dx.
\]
These are exactly the moments of~\eqref{aud:fr:measure}; the zeroth moment is $k/d\to\kappa$. Polynomial approximation on the common compact support proves weak convergence. The case $\kappa=0$ also follows directly from the vanishing total mass. Uniformity follows by the same estimates, or by extracting a convergent subsequence from any hypothetical sequence violating it.
\end{proof}

Lemma~\ref{aud:fr:measurelemma} is an elementary specialization of the classical varying-recurrence zero-distribution theorem of Kuijlaars and Van Assche~\cite[Theorem 1.4]{kuijlaars}. The proof above is included so that no general zero-distribution theorem is required for the argument.

\section{The exterior saddle and its regular branch}
Define
\begin{equation}\label{aud:fr:Laurent}
 D(v)=A(v)B(v),\qquad
 A(v)=\prod_{j=1}^m(1+v/j^2),\qquad
 B(v)=\prod_{\lambda>0}(1-\alpha/v).
\end{equation}
For $\rho>\max\alpha$, the value $D(\rho)$ is positive, and
\begin{align}
 \mu(\rho)&=\frac{\rho D'(\rho)}{D(\rho)}
 =\sum_{j=1}^m\frac{\rho}{\rho+j^2}
   +\sum_{\lambda>0}\frac{\alpha}{\rho-\alpha},
 \label{aud:fr:mean}\\
 V(\rho)&=\frac{d\mu}{d\log\rho}=V_A-V_B,
 \label{aud:fr:curvaturefinite}\\
 V_A&=\sum_{j=1}^m\frac{\rho j^2}{(\rho+j^2)^2},\qquad
 V_B=\sum_{\lambda>0}\frac{\alpha\rho}{(\rho-\alpha)^2}.
 \nonumber
\end{align}
Write $\rho=d^2y$. If $k/d\to\kappa$ and $y>\kappa(\kappa+4)$, the norm bound separates this radius from all roots. Lemma~\ref{aud:fr:measurelemma} gives
\begin{equation}\label{aud:fr:meanlimit}
 \frac{\mu(d^2y)}d\longrightarrow \mathcal M(\kappa,y)
 =\sqrt y\arctan\frac2{\sqrt y}
 +\frac12\int_0^\kappa
 \left(\sqrt{\frac y{y-x(x+4)}}-1\right)\,dx.
\end{equation}
The factor $1/2$ counts one root from each symmetric pair; a possible zero root contributes zero. The convergence, including derivatives in $y$, is uniform on compact subsets of the strict exterior region. Indeed all the rational test functions and their derivatives are uniformly bounded there.

Put
\begin{equation}\label{aud:fr:J}
 \mathcal J(\kappa,y)=\int_0^\kappa\frac{dx}{\sqrt{y-x(x+4)}}
 =\arcsin\frac{\kappa+2}{\sqrt{y+4}}
  -\arcsin\frac2{\sqrt{y+4}}.
\end{equation}
Since $h/d\to(1-\kappa)/2$, the limiting equation $\mu(\rho)=h$ becomes
\begin{equation}\label{aud:fr:saddleequation}
 \sqrt y\left(\arctan\frac2{\sqrt y}
       +\arcsin\frac{\kappa+2}{\sqrt{y+4}}\right)=1.
\end{equation}
Equivalently, $2\sqrt y\arctan(2/\sqrt y)+\sqrt y\,\mathcal J(\kappa,y)=1$.

To describe the relevant branch explicitly, define
\begin{align}
 K(t)&=\frac{\sin t}{t}-2\cos t-2,
 \label{aud:fr:K}\\
 F(t)&=(2t^2-1)\sin t+t\cos t,\qquad
 D_0(t)=\cos t+2t\sin t.
 \label{aud:fr:FD}
\end{align}
Let $t_0$ be the unique nonzero solution of $\tan(t_0/2)=2t_0$ in $(0,\pi)$. The derivative of $\tan(t/2)-2t$ increases strictly on this interval, begins negative, and tends to infinity; this proves uniqueness. Evaluation at $3\pi/4$ gives $t_0>3\pi/4$. The earlier fixed-degree parameter satisfies $t_0=1/(2a_0)$, where $a_0\arctan(1/a_0)=1/4$.

\mergenote{The ``earlier fixed-degree parameter'' is the constant $a$ of manuscript~64, written $a_0$ in~\eqref{aud:eq:wa}. Manuscript~58 does not cite~64. That manuscript's fixed-degree phase laws are printed as the case $\kappa=0$ in Sections~\ref{aud:sec:kappa0} and~\ref{aud:sec:kappa0k}.}

\begin{proposition}[The positive-curvature branch]\label{aud:fr:branch}
There is a unique $t_c\in(t_0,\pi)$ with $F(t_c)=0$. Set
\[
 \kappa_c=K(t_c).
\]
The function $K$ increases strictly from zero to $\kappa_c$ on $[t_0,t_c]$. For every $0\le\kappa<\kappa_c$, let $t=t(\kappa)$ be its inverse on this interval. Then
\begin{equation}\label{aud:fr:parametric}
 y(\kappa)=\frac1{t(\kappa)^2}
\end{equation}
solves~\eqref{aud:fr:saddleequation}, lies strictly outside the root support, and has limiting curvature
\begin{equation}\label{aud:fr:curvaturelimit}
 yM_y(\kappa,y)=\frac{F(t)}{4tD_0(t)}>0.
\end{equation}
At $\kappa_c$ the curvature is zero while the root gap is still positive. No global uniqueness assertion for all exterior saddles is made.
\end{proposition}
\begin{proof}
One has $K(t_0)=0$, $K'(t)=F(t)/t^2$, and
\[
 F(t_0)=\frac{t_0(4t_0^2-3)}{1+4t_0^2}>0,\qquad F(\pi)=-\pi.
\]
For $3\pi/4<t<\pi$,
\[
 F'(t)=t(3\sin t+2t\cos t)<0,
\]
because $\sin t\le-\cos t$ and $2t>3$. Thus $F$ has exactly one zero in $(t_0,\pi)$, proving the asserted monotonicity of $K$.

On $[t_0,t_c]$, $D_0$ is positive: $D_0(t_0)=1$, its derivative $\sin t+2t\cos t$ is negative, and $F(t_c)=0$ gives $D_0(t_c)=\sin(t_c)/t_c>0$. Direct substitution into~\eqref{aud:fr:K} gives the exact gap identity
\begin{equation}\label{aud:fr:gap}
 \frac1{t^2}-K(t)(K(t)+4)
 =\frac{D_0(t)^2}{t^2}>0.
\end{equation}
Furthermore, $t-\arctan(2t)$ is positive and less than $\pi/2$: it lies in $(0,\pi)$, and its cosine has the positive sign of $D_0(t)$. Therefore
\[
 \arcsin\frac{K(t)+2}{\sqrt{t^{-2}+4}}
 =t-\arctan(2t),
\]
with the principal inverse-sine branch. This proves~\eqref{aud:fr:saddleequation}.

Finally set $E(\kappa,y)=2\mathcal M(\kappa,y)+\kappa-1$. From~\eqref{aud:fr:saddleequation},
\[
 E_\kappa(\kappa,y)=\frac{\sqrt y}{\sqrt{y-\kappa(\kappa+4)}}.
\]
Differentiate $E(K(t),t^{-2})=0$ and use~\eqref{aud:fr:gap}. The result is
\[
 yM_y=\frac{K'(t)}{4\sqrt{y-\kappa(\kappa+4)}}
      =\frac{F(t)}{4tD_0(t)}.
\]
It is positive before $t_c$ and zero there, whereas~\eqref{aud:fr:gap} remains strictly positive at $t_c$.
\end{proof}

\mergenote{Shared statement. Proposition~\ref{aud:fr:branch}, the gap identity~\eqref{aud:fr:gap} and the curvature~\eqref{aud:fr:curvaturelimit} are restated, with their elementary proofs, in manuscripts~49 (its Section~2), 53 (Section~2), 54 (Section~3) and~40 (Section~3, in the variable $T$). Manuscripts~34 and~37 use the point $T=14/5$ of the same branch. These facts are printed once, here, and the members' sections refer to them.}

For orientation only, the implicit constants are approximately
\[
 t_0\approx2.78649815065118,\qquad
 t_c\approx2.96463500776812,\qquad
 \kappa_c\approx0.0281461097171248.
\]
These decimals are not used to prove the branch statements.

\begin{lemma}[An explicit rational subinterval]\label{aud:fr:rationalband}
One has $\kappa_c>1/36$.
\end{lemma}
\begin{proof}
At $x=74/25$, let
\[
 S_n(x)=\sum_{j=0}^n\frac{(-1)^jx^{2j+1}}{(2j+1)!},\qquad
 C_n(x)=\sum_{j=0}^n\frac{(-1)^jx^{2j}}{(2j)!}.
\]
The alternating Taylor bounds give $S_5<\sin x<S_6$ and $C_5<\cos x<C_6$. Exact rational arithmetic gives
\[
 \frac{S_5(x)}x-2C_6(x)-2>\frac1{36},\qquad
 (2x^2-1)S_5(x)+xC_5(x)>0.
\]
The inequalities can also be verified directly by clearing factorials and powers of $25$; an accompanying exact-rational script records both positive differences. The elementary bounds $31/10<\pi<22/7$ place $x$ in $(3\pi/4,\pi)$. Hence $K(x)>1/36$ and $F(x)>0$, so Proposition~\ref{aud:fr:branch} gives $t_0<x<t_c$ and $\kappa_c>1/36$.
\end{proof}

\mergenote{Manuscripts~49 and~53 delivered the same exact-rational certificate, with the same script and output, byte for byte. They use it to place their base intervals $[1/200,1/100]$ and $[1/100,1/50]$ inside $(0,\kappa_c)$ (Sections~\ref{aud:sec:qgbase} and~\ref{aud:sec:53a}). It is shipped once, as \nolinkurl{code/58-finite-ratio-verify_rational_band.py} with its record \nolinkurl{data/58-finite-ratio-rational_band_certificate.json}.}

\section{Uniform positive saddle coefficients}
\begin{lemma}\label{aud:fr:Gaussian}
Fix a compact interval $I\subset[0,\kappa_c)$. For all sufficiently large $d$ and integers $k$ with $k/d\in I$, there is a solution $\rho$ of $\mu(\rho)=h$ in a neighborhood of $d^2y(k/d)$, unique in that neighborhood. Uniformly over these pairs,
\begin{equation}\label{aud:fr:finitesaddle}
 \frac{\rho}{d^2}-y(k/d)\longrightarrow0,\qquad
 \frac{V(\rho)}d-\frac{F(t(k/d))}{4t(k/d)D_0(t(k/d))}\longrightarrow0.
\end{equation}
Let $C_h=[v^h]D(v)$. Then
\begin{equation}\label{aud:fr:positivecoefficient}
 C_h=\frac{D(\rho)\rho^{-h}}{\sqrt{2\pi V(\rho)}}(1+o(1))>0.
\end{equation}
For every fixed integer $j$,
\begin{equation}\label{aud:fr:coefficientratio}
 \frac{[v^{h-j}]D(v)}{\rho^j C_h}\longrightarrow1
\end{equation}
uniformly on the same interval of ratios.
\end{lemma}
\begin{proof}
The limiting root gap and limiting curvature have positive uniform lower bounds on $I$. The norm estimate~\eqref{aud:fr:rootbound} and uniform convergence in~\eqref{aud:fr:meanlimit}, including its derivative, therefore provide a common neighborhood of the branch on which the finite mean is strictly increasing and its values surround $h$. The intermediate value theorem gives the saddle and~\eqref{aud:fr:finitesaddle}.

Write $x_\lambda=\alpha/\rho$. For every real $\theta$ there is the exact bound
\begin{equation}\label{aud:fr:angular}
 \left|\frac{D(\rho e^{i\theta})}{D(\rho)}\right|
 \le\exp\bigl(-2(V_A-V_B)\sin^2(\theta/2)\bigr).
\end{equation}
Indeed, for an $A$ factor the squared normalized modulus is $1-4p(1-p)\sin^2(\theta/2)$, with $p=\rho/(\rho+j^2)$. For a $B$ factor it is
\[
 \left|\frac{1-x_\lambda e^{-i\theta}}{1-x_\lambda}\right|^2
 =1+\frac{4x_\lambda\sin^2(\theta/2)}{(1-x_\lambda)^2}.
\]
Apply $\log(1+z)\le z$ to both products. This proves~\eqref{aud:fr:angular}; if an $A$ factor vanishes, the desired bound follows immediately.

The quantities $x_\lambda$ are uniformly bounded away from one. There are $O(d)$ factors, and every third logarithmic derivative near $\theta=0$ has a common bound after summation. Consequently
\[
 \log\frac{D(\rho e^{i\theta})}{D(\rho)}-ih\theta
 =-\frac{V(\rho)}2\theta^2+O(d|\theta|^3)
\]
uniformly. On $|\theta|\le d^{-2/5}$ the remainder tends to zero; outside this arc,~\eqref{aud:fr:angular} and $V(\rho)\ge c_I d$ give exponential decay. Cauchy's coefficient formula, valid for Laurent polynomials on every circle of positive radius, now gives~\eqref{aud:fr:positivecoefficient} by the Gaussian integral. Inserting $e^{ij\theta}$ gives~\eqref{aud:fr:coefficientratio}. The leading expression in~\eqref{aud:fr:positivecoefficient} is positive, proving eventual positivity of $C_h$.
\end{proof}
